% ============================================================================
% ch11-e.tex —— §11.5.2 阻抗技术在电催化领域中的应用（前半）
%               现场/外场 EIS 的取舍、参比与工作电极制备、两种测量模式
% 源：sources/parts/ch11-b.txt 行 1095 起；书内 p.862—865 / PDF p.350—353
%
% 【节序说明 —— 装配者必读】
%   本文件在 chapters/ch11.tex 的 \input 顺序中排在最后（a→b→d→c→e），
%   但按源书节序，§11.5.2 应排在 §11.5.1（对应 part d）之后、
%   §11.6（对应 part c 后半）之前。即正确阅读顺序是 a→b→d→e→c，
%   而不是装配文件当前的 a→b→d→c→e。
%   本文件**不**修改 chapters/ch11.tex，节序调整由装配者统一处理。
% ============================================================================

\subsection{阻抗技术在电催化领域中的应用}

\src{书内 p.862 / PDF p.350}

前面讨论的都是同一个电极过程内部的等效元件。一旦把 EIS 用到一个真实的电催化
体系上，问题就分裂成两个层面：\emph{在哪里测}（现场还是外场）与\emph{怎么测}
（何种电极、何种扰动模式）。这两个层面的选择往往比后续的等效电路拟合更能
决定数据是否可用，因此本节先讨论它们，再进入典型谱图的判读。

\subsubsection{现场与外场阻抗谱}

\begin{principle}
\textbf{按测量对象划界的两种 EIS}。取决于研究目的，EIS 可以分为现场的
（in-situ）与外场的（ex-situ）。这一划分不是按仪器或频段，而是按\emph{被测
体系在测量时是否处于其真实工作状态}：现场 EIS 把单池或电堆整体作为被测
对象，在电池实际运行、带直流负载的条件下叠加交流扰动；外场 EIS 则把材料或
组件（催化剂、催化层、膜、GDE）取出，放进一个可控的电解池里单独表征。分界
线一旦划定，后面所有的取舍——装置成本、可获得的参数、数据的可解释性——
都由它派生出来。
\end{principle}

\begin{measure}
\textbf{现场 EIS：装置、代价与人为因素}。现场 EIS 经常用于单池或电堆的分析。
它的测量装置通常是复杂、昂贵且耗时的。更棘手的是运行期间——尤其在高直流
负载时——由于电缆和连线等的干扰，可能引入人为因素，使谱的高频段不再反映
电池本身。这里的因果关系值得说清楚：高直流负载意味着大电流，大电流在电缆
与连线上产生可观的电压降与感抗，而感抗项随频率上升而增大，于是高频端出现的
那个“电感尾巴”究竟来自电池还是来自测量回路，无法从谱形本身分辨。这也是
现场 EIS 最典型的失效方式。
\end{measure}

\begin{measure}
\textbf{外场 EIS：三电极体系与电解质选择}。正因为上述代价，外场 EIS 一般是
很多燃料电池研究者的首选。外场 EIS 一般采用传统的三电极体系在硫酸电解质中
进行，参比电极一般为饱和甘汞电极（SCE）、标准氢电极（NHE）、\ce{Ag/AgCl}
或 \ce{Hg/HgO} 电极，有时也用动态氢电极（DHE）作参比电极；工作电极一般通过
将薄层催化剂铸在干净的玻碳电极或热解石墨电极上，或者将电极直接浸泡在含有
欲测催化剂的溶液中，有时，气体扩散电极（GDE）也可以用作工作电极。
\src{书内 p.863 / PDF p.351}
\end{measure}

\begin{probout}
\textbf{现场与外场 EIS 的取舍判据}——不是优劣排序，而是“什么目的配什么代价”：

\begin{itemize}
  \item \textbf{要问“这个电池/电堆现在出了什么问题”}$\rightarrow$ 选现场 EIS。
        只有现场测量才能保留膜、气体扩散电极、MEA、单池各元件在真实工作
        条件下的贡献，也才可能定位问题元件。\emph{代价}：装置复杂、昂贵、
        耗时；高直流负载下电缆与连线引入人为因素，高频段不可信。
  \item \textbf{要问“这个催化剂/材料本身好不好”}$\rightarrow$ 选外场 EIS。
        三电极体系把工作电极电位测准，排除了对电极与欧姆降的纠缠，数据
        可解释性最高。\emph{代价}：离开真实工况，电解质、温度、气体分压、
        水管理都与单池不同，测得的活性排序未必能外推到电堆。
  \item \textbf{两者都要}时，正确用法是外场 EIS 建立机理与参数基准，现场 EIS
        检验该基准在工况下是否仍然成立，而不是把两者的数值直接相减当作
        “工况劣化量”。
\end{itemize}
\end{probout}

\subsubsection{电极制备与测量模式}

\begin{dataflow}
\textbf{从电极制备到一条可信的阻抗谱}。参数的可信度在电流进入恒电位仪之前
就已经被决定了，链条如下：

\begin{enumerate}
  \item \textbf{参比电极选型}（SCE / NHE / \ce{Ag/AgCl} / \ce{Hg/HgO} / DHE）
        $\rightarrow$ 决定电位标度与液接电位误差；
  \item \textbf{工作电极制备}（薄层催化剂铸于干净玻碳或热解石墨，或电极直接
        浸泡于含催化剂的溶液，或用 GDE）$\rightarrow$ 决定催化剂的真实负载量、
        分散状态与接触电阻；
  \item \textbf{扰动模式选择}（恒电位 / 恒电流）与交流振幅、频率范围
        $\rightarrow$ 决定谱是否在仪器的线性区内被测得；
  \item \textbf{等效电路拟合} $\rightarrow$ 输出 $R_\mathrm{s}$、
        $C_\mathrm{dl}$、$R_\mathrm{ct}$ 等参数。
\end{enumerate}

这条链条是\emph{串联}的：第 1、2 环节引入的误差无法在第 4 环节通过拟合消除，
只会在拟合结果里伪装成某个元件参数。
\end{dataflow}

\begin{pitfall}
\textbf{为什么参比与工作电极制备决定数据可信度}。三点具体的因果：

\begin{itemize}
  \item \textbf{参比电极决定“电位”这个自变量的真值。} 不同参比（SCE、NHE、
        \ce{Ag/AgCl}、\ce{Hg/HgO}、DHE）各自的电极电位不同，报告数据时必须
        写明参比与电解质；只写“在 0.5 V 下测得”而不写参比，等于没有给出
        电位。DHE 之所以称为“动态”氢电极，正是因为它的电位随体系中氢分压
        与酸度而移动，用它可以减小与实际工况的偏离，但也就失去了固定标度的
        便利。
  \item \textbf{工作电极制备决定“有多少催化剂在工作”。} 薄层催化剂铸在玻碳
        或热解石墨上，要求基底干净：基底上的残留物既贡献额外的界面电容，
        又可能自身有电化学活性。而“将电极直接浸泡在含有欲测催化剂的溶液中”
        这一做法，催化剂是悬浮或溶解状态，其负载量与附着状态无法定量，得到
        的阻抗只能作相对比较。GDE 则引入了额外的气体扩散层与三相界面，谱中
        会多出传质相关的特征。
  \item \textbf{两种制备方式的差异会在谱上表现为同一个 $R_\mathrm{ct}$ 的不同
        数值}，而谱形本身无法告诉你究竟是本征活性不同，还是接触与负载不同。
        这正是必须在论文中完整描述电极制备细节的原因。
\end{itemize}
\end{pitfall}

\begin{measure}
\textbf{两种测量模式及其各自的代价}。阻抗测量一般有两种模式，一种是恒电位
模式，另一种是恒电流模式。从这两种模式中得到的结果基本上没有明显不同。然而，
取决于测量对象，这两种模式各有优缺点。\src{书内 p.862 / PDF p.350}

\emph{恒电位模式}：扰动是交流电位，响应是通过研究体系的测量电流。由于一个
小的电位都可能导致大的电流变化，从而导致恒电位仪的过载，因此，尝试不同的
交流电压确保恒电位仪不会过载是非常必要的。在燃料电池的测量中，交流振幅一般
选择在 5$\sim$15 mV，扫描频率范围一般在 1 mHz$\sim$100 kHz 之间。在催化剂
活性的测量中，控制电位能更好地控制电化学反应的发生，因此一般采用恒电位模式。

\emph{恒电流模式}：扰动是交流电流，响应是通过研究体系的测量电位。由于交流
和直流大小很容易被运算放大器所确定，能有效防止电池和恒电位仪过载，因此能
采用相对大的交流电，从而得到更多的阻抗信息。对于电池阻抗测量，一般电流比较
大，因此更适合用恒电流模式测量，恒电位模式容易导致过载。
\src{书内 p.862 / PDF p.350}
\end{measure}

\begin{center}
\small
% 注意：本项目 calc 包会改写 \dimexpr，任何 `p{\dimexpr...\onecolwidth\relax}`
% 形式的尺寸运算都会报 "Illegal unit of measure"（calc Error: `.' invalid）。
% 故首列改用固定 26 mm（版心 165 mm），其余由 tabularx 按 \onecolwidth 分配，
% 全程不做尺寸运算。\onecolwidth = 0.80\textwidth ≈ 132 mm。
\begin{tabularx}{\onecolwidth}{@{}>{\raggedright\arraybackslash}p{26mm}
  >{\raggedright\arraybackslash}X >{\raggedright\arraybackslash}X@{}}
\toprule
 & 恒电位模式 & 恒电流模式 \\
\midrule
扰动量 & 交流电位 & 交流电流 \\
响应量 & 测量电流 & 测量电位 \\
主要风险 & 小电位变化即可引起大电流，恒电位仪过载 &
            交流与直流幅值不易独立设定 \\
适用对象 & 催化剂活性测量（控制电位即控制反应） &
            电池与电堆阻抗（电流大，恒电位易过载） \\
常用参数 & 燃料电池中振幅 5$\sim$15 mV；频率 1 mHz$\sim$100 kHz &
            可采用相对大的交流电，阻抗信息更多 \\
\bottomrule
\end{tabularx}
\end{center}

\begin{pitfall}
\textbf{四类会把谱读错的误判}。

\begin{enumerate}
  \item \textbf{把接触电阻当电荷传递电阻。} 催化剂层与玻碳或热解石墨基底之间
        的接触电阻、以及电缆与接头引入的欧姆项，都会在 Nyquist 图上表现为
        高频端的实轴截距或一个额外的小半圆。若不先做空白基底测量作对照，
        很容易把这个本不属于电极反应的电阻读成 $R_\mathrm{ct}$，从而低估
        催化活性。判据是：$R_\mathrm{ct}$ 随电位（过电位）应有系统性变化，
        接触电阻不会。
  \item \textbf{恒电位仪过载导致谱失真。} 恒电位模式下，小的电位扰动即可引起
        大电流，一旦超出恒电位仪的电流量程，仪器不再维持设定的交流电位，谱形
        被削平或出现伪感抗，而软件仍会照常输出数据点。因此“尝试不同的交流
        电压确保恒电位仪不会过载是非常必要的”——这句话不是操作建议，而是
        数据有效性的前提。
  \item \textbf{外场 EIS 的电解质与实际工况不符。} 外场 EIS 常在硫酸电解质中
        用三电极体系进行，而真实燃料电池是固体聚合物电解质、含液态水与气体
        扩散层。电解质种类、离子强度、pH 与传质环境都不同，因此在硫酸中测得
        的 $R_\mathrm{ct}$ 排序不能直接当作电堆中的性能排序。
  \item \textbf{现场 EIS 的电缆干扰。} 高直流负载下，大电流经电缆与连线产生的
        干扰会把测量回路自身的阻抗叠加到谱上，通常表现为高频区的伪感抗或
        不可解释的相角。判据是：更换电缆走线或缩短线长后高频段发生变化，
        则该段特征是人为因素而非电池特征。
\end{enumerate}
\end{pitfall}

\begin{crosscheck}
\textbf{读图实例：原书图 11-41 的 Nyquist 图}。该图为溶液电阻不可忽略时电化学
极化与浓差极化同时存在电极的等效电路及 Nyquist 图。
\src{书内 p.862 / PDF p.350}

先说几何，再说归属。对于电化学极化和浓差极化同时存在的电极，该电极过程由
电荷传递过程和扩散过程共同控制。
\end{crosscheck}

\begin{obs}
谱形由两段构成：

\begin{itemize}
  \item 高频区：一个半圆，其半径为 $R_\mathrm{ct}/2$；
  \item 低频区：一条 45° 的直线。
\end{itemize}
\end{obs}

\begin{infer}
由此可读出控制步骤的归属：

\begin{itemize}
  \item 半圆表示电荷传递过程为控制步骤——半圆的直径给出 $R_\mathrm{ct}$；
  \item 45° 直线表示由电极反应的反应物或产物的扩散控制；
  \item 因此该电极过程由电荷传递过程和扩散过程\emph{共同}控制。
\end{itemize}
\end{infer}

\begin{crosscheck}
\textbf{观察与推断的界线必须守住}。图谱本身只给出“半圆 + 45° 直线”这一几何
事实；半圆 $\rightarrow$ 电荷传递控制、直线 $\rightarrow$ 扩散控制，是建立在
等效电路模型之上的解释。若换了模型（例如存在两个时间常数、半圆被压扁、或
直线偏离 45°），同一段几何形状可以有别的解释。
\end{crosscheck}

\begin{pitfall}
读原书图 11-41 时最容易犯的两个错：一是把高频端半圆在实轴上的\emph{截距}
当成 $R_\mathrm{ct}$——只有当溶液电阻可忽略时截距才近似等于
$R_\mathrm{ct}$，而此图的前提恰恰是“溶液电阻不可忽略”，所以必须先扣除溶液
电阻 $R_\mathrm{s}$，再取半圆直径；二是把低频 45° 直线的起点当作扩散开始
出现的位置精确定值，实际上直线段与半圆的过渡是渐变的，起点的判定依赖拟合
区间的人为选择。
\end{pitfall}

此处已对照原页面图像（书内 p.862 / PDF p.350）：原书图~11-41 确为高频区一个
半径 $R_\mathrm{ct}/2$ 的半圆加低频区一条 $45^\circ$ 的直线，图中另以竖直虚线标出
“电荷传递控制”与“扩散控制”的分界；等效电路元件印作 $R_\Omega$、$R_\mathrm{ct}$、
$C_\mathrm{dl}$ 与方框化的 $\tilde Z_\mathrm{W}$，半圆一侧的实轴标出 $R_\Omega$、
$R_\Omega+R_\mathrm{ct}/2$、$R_\Omega+R_\mathrm{ct}$。与本节的谱形叙述一致。

此处已对照原页面图像（书内 p.861--862 / PDF p.349--350）：图~11-40（p.861）与
图~11-41（p.862）的图题\textbf{印刷完全相同}，均为“溶液电阻不可忽略电化学极化
电极的等效电路及 Nyquist 图”，属原书自身图题重复；正文分别以“其等效电路及
Nyquist 如图 11-40 所示”“……如图 11-41 所示”引用两图，本节按正文分工处理无误。

此处已对照原页面图像（书内 p.862 / PDF p.350）：原书正文印“交流振幅一般选择在
5$\sim$15 mV，扫描频率范围一般在 1 mHz$\sim$100 kHz 之间”，数值上下限与本处一致
（原书用全角波浪号“～”，此处用 $\sim$，属排版差异）。

此处已对照原页面图像（书内 p.861--862 / PDF p.349--350）：原书印“参比电极一般为
饱和甘汞电极（SCE）、标准氢电极（NHE）、Ag/AgCl 或 Hg/HgO 电极，有时也用动态氢
电极（DHE）作参比电极”，\ce{Hg/HgO} 的写法与“动态氢电极（DHE）”的中文名、缩写
均无误。

\src{书内 p.862--863 / PDF p.350--351}
